\section{Assembly of the rational trace formula}
\label{sec:rational-assembly}
We now combine the operator estimates, keeping the principal polar
coordinates separate until the final trace.

\begin{proof}[Proof of Theorem~\ref{thm:main-dirichlet}]
Lemma~\ref{lem:raw-estimates}, Lemma~\ref{lem:theta-identities},
and Proposition~\ref{prop:gram-control} establish the concentration,
graph-domain, Poisson, and inverse-Gram inputs of
Theorem~\ref{thm:rational-local-limit}. The principal evaluation
repair is part of that theorem's proof. Thus its conclusion is
unconditional for every fixed primitive character.
Theorem~\ref{thm:unconditional-frequency-tightness} upgrades this
local convergence to Schwartz frequency tests. Since $h$ is
compactly supported and smooth, $\widehat h(\tau)=H(i\tau)$ is
Schwartz. The exact Fourier-diagonal formula is
\[
 \int\widehat h(\tau)d\mu_\Lambda(\tau)
   =2Lh(0)-\Tr(S_LQ_{\chi,\Lambda}S_LV(h)).
\]
One may define the trace on the left-supported interval by
compressing $V(h)$ on both sides; its smooth kernel on a compact
interval is trace class. The same equality follows first for
convolution squares from the Hilbert--Schmidt distance formula
and then by polarization. The preceding convergence gives the
limit $\mathfrak B_\chi(h)$.
Proposition~\ref{prop:trace-compression} replaces the compressed
projection by $Q_{\chi,\Lambda}$ and proves
\eqref{eq:main-sector-limit}. In the principal sector
Proposition~\ref{prop:polar-trace} gives exactly the two scalar
characters in \eqref{eq:main-principal-limit}. No step in this
chain uses a theorem placing nontrivial zeros on the line.
\end{proof}

\begin{proof}[Proof of Corollary~\ref{cor:main-adelic}]
Haar orthogonality on $\widehat\Z^\times$ selects the coefficient
$h_\chi$ in each character block. There are finitely many such
blocks, and $h(1,1)=\sum_\chi h_\chi(0)$. Sum
Theorem~\ref{thm:main-dirichlet} over them. The polar coordinates
occur only in the principal block and cancel the explicit polar
integral in $T_\Lambda$.

If all zeros are on the line, each harmonic measure is a Dirac
mass, so the result equals the zero side of Weil's explicit
formula. Conversely, equality with Weil's distribution for all
Bruhat tests makes that distribution nonnegative on convolution
squares, because each $\mathfrak B_\chi$ is a positive measure.
Weil's criterion, in the precise global form of
\citet[Proposition~6.2 and Corollary~6.3]{ccm2007}, then puts all
nontrivial zeros in all sectors on the critical line. Finite
linear combinations of characters allow the sectors to be
tested separately. This proves the equivalence.
Using Alp\"oge's Dirichlet theorem, with the externally verified
scope described in \citet{mitverification}, gives its stated
trace-formula consequence.
\end{proof}

The sign argument does not order $\mathfrak B$ and $\mathfrak W$.
It uses equality of two distributions, one of which is already
positive. Proposition~\ref{prop:sign} explains why replacing
this step by a definite sign for their difference would be invalid.

Only two consequences of the prolate transition distribution
enter this proof: a positive lower bound on the selected
concentration eigenvalues and an upper bound below one on the
first omitted eigenvalue. The dimension-free smoothed-tail and
band-extension estimates avoid summing the individual leakage
defects across the transition. No transition indices are removed
from the cutoff in Theorem~\ref{thm:main-dirichlet}.
